```latex
\chapter{Conclusiones y Recomendaciones}
\label{chap:conclusiones}

\section{Conclusiones}

\subsection{Validación del Objetivo General y la Hipótesis Ciberfísica}
En el contexto de la presente investigación de maestría, hemos logrado demostrar empíricamente la viabilidad técnica y la validez funcional de implementar un Gemelo Digital de Nivel 3 con capacidades predictivas e interactivas sobre una máquina dosificadora de granos de naturaleza \textit{legacy}. La hipótesis fundamental del proyecto sostenía que era posible modernizar un equipo industrial obsoleto sin alterar su arquitectura electromecánica original mediante un \textit{bypass} ciberfísico no intrusivo. Los resultados cuantitativos obtenidos durante las jornadas de validación experimental en el Laboratorio de Procesos confirman esta premisa. La integración efectiva de los tres pilares del sistema—el \textit{hardware} de instrumentación externa basado en controladores ESP32, el \textit{middleware} de comunicación soportado por el protocolo MQTT y el broker Mosquitto, y el entorno de simulación 3D en Webots—constituye un sistema ciberfísico cohesivo.

La evidencia más contundente del cumplimiento del objetivo general reside en la consecución de una Efectividad General de los Equipos (OEE) promedio del 91.85\% con una desviación estándar de apenas ±0.46\% durante los ciclos operativos normales. Esta métrica, complementada por una Disponibilidad del 100\% y un Rendimiento equivalente fuera de los escenarios de fallo, demuestra que el gemelo digital no solo funciona como un espejo pasivo del proceso físico, sino que mantiene la estabilidad y eficiencia del sistema. La capacidad predictiva, característica distintiva del Nivel 3 en la pirámide de los gemelos digitales, fue validada de manera crítica en la Prueba 5. En dicho escenario, se inyectó un fallo en el sistema neumático que fue detectado y reflejado por el modelo virtual, permitiendo una reacción del operador asistido que resultó en un Tiempo Medio de Reparación (MTTR) de 52 segundos y una disponibilidad durante el evento de fallo del 92.92\%. Estos valores confirman que el gemelo digital trasciende la mera monitorización para convertirse en una herramienta de apoyo a la decisión en tiempo real.

\subsection{Cumplimiento de los Objetivos Específicos y Ejecución Metodológica}
El desarrollo del proyecto fue guiado por un marco metodológico estructurado que se desglosó en tres objetivos específicos, cada uno de los cuales fue alcanzado a cabalidad y contribuyó de manera sinérgica al éxito del prototipo final.

En primer lugar, el objetivo específico OE1 (Levantamiento y diagnóstico) se cumplió mediante un proceso exhaustivo de ingeniería inversa. Ante la falta de documentación técnica actualizada del equipo \textit{legacy}, se realizó un mapeo detallado del circuito eléctrico a 220V y de la lógica de control original gobernada por un PLC. Este diagnóstico permitió identificar el punto exacto de intervención en la entrada 1 del PLC, proveniente del contactor trifásico, lo cual fue fundamental para diseñar el \textit{bypass} sin desencadenar conflictos lógicos ni riesgos eléctricos. La instalación de la sensórica (infrarrojos, ultrasónico y transductor de presión) respondió directamente a la necesidad de caracterizar el comportamiento neumático y mecánico previamente oculto.

En segundo lugar, el objetivo OE2 (Diseño e implementación) se concretó con la materialización completa de la arquitectura propuesta. Kiany Cárdenas lideró el diseño del \textit{hardware} de la caja de control externa, resolviendo las limitaciones de pines de los ESP32 mediante módulos de expansión I2C y mitigando las interferencias electromagnéticas con un condensador de poliéster dimensionado para la carga inductiva de los contactores. Paralelamente, Bryan Toaza desarrolló la arquitectura de \textit{software}, implementando la Máquina de Estados Finitos (FSM) en Python dentro de Webots y programando los nodos de comunicación y el SCADA en Node-RED. La convergencia de ambos desarrollos resultó en una latencia de comunicación MQTT promedio de 0.61 ms, un valor excepcionalmente por debajo del umbral de 100 ms establecido en los requerimientos no funcionales, garantizando una sincronía efectiva entre el mundo físico y el virtual.

Finalmente, el objetivo OE3 (Evaluación) fue acometido mediante un riguroso plan de pruebas que procesó lotes de producción desde 60 hasta 200 fundas. Esta fase no se limitó a verificar la funcionalidad, sino que generó un conjunto de datos cuantitativos robustos. La repetibilidad de la métrica OEE en las cuatro pruebas sin fallo, con una variación mínima, valida la consistencia de la instrumentación y la robustez del algoritmo de control. La ejecución completa del proyecto dentro de un presupuesto de $243.20 USD valida que los objetivos técnicos se alcanzaron bajo estrictas restricciones económicas, maximizando la relación costo-beneficio.

\subsection{Aportaciones Técnicas al Campo de los Sistemas Ciberfísicos y la Manufactura}
Más allá de la solución puntual, este trabajo de maestría deja un legado de aportaciones técnicas y metodológicas relevantes para la Ingeniería en Tecnologías de la Información. La principal contribución es la formalización de una metodología de \textit{bypass} ciberfísico no intrusivo para la modernización de máquinas \textit{legacy}. A diferencia de un \textit{retrofit} tradicional que requiere la modificación completa del tablero de control, nuestro enfoque permite instrumentar externamente el proceso, preservando la integridad del circuito original de 220V y la lógica del PLC. Esta metodología minimiza el riesgo de inutilizar equipos críticos cuya detención prolongada no es viable.

En segundo lugar, se ha validado una arquitectura de integración vertical de bajo costo que desafía el paradigma de que los gemelos digitales requieren plataformas comerciales de alto valor. Utilizando microcontroladores ESP32 de amplia disponibilidad, protocolos ligeros como MQTT y plataformas de simulación de código abierto como Webots, se logró un sistema plenamente funcional descrito en detalle teórico y práctico. Este proyecto se constituye como un caso de estudio replicable y totalmente documentado para entornos académicos y de formación técnica, proporcionando un manual para la transformación digital de laboratorios universitarios. En el ámbito de la interoperabilidad, se ha demostrado cómo un protocolo IoT moderno puede inyectar datos de un activo \textit{legacy} en un ecosistema digital, un paso fundamental para la Industria 4.0.

\subsection{Identificación de Limitaciones y Restricciones del Prototipo Actual}
Es imperativo para el rigor científico reconocer las limitaciones que enmarcaron el alcance de este proyecto. La restricción más significativa fue la imposibilidad de medir la Calidad dentro del indicador OEE. La ausencia de un sistema de pesaje dinámico (celdas de carga) o de un subsistema de visión artificial en la etapa de dosificación impidió cuantificar la precisión en gramos de cada funda. Por lo tanto, el OEE reportado (91.85\%) es una métrica compuesta únicamente por Disponibilidad y Rendimiento, lo cual subestima potencialmente las pérdidas del proceso si existiesen variaciones no detectadas en la masa del producto. Esta realidad fue una consecuencia directa de la naturaleza \textit{legacy} de la máquina y de la falta de permisos para modificar su tolva de pesaje original.

En segundo lugar, el dominio de validación estuvo circunscrito a un entorno controlado de laboratorio. Si bien esto es adecuado para una prueba de concepto a nivel de maestría, las condiciones no replican completamente las perturbaciones ambientales de una planta agroindustrial real, como el polvo en suspensión, las fluctuaciones extremas de temperatura o las vibraciones constantes transmitidas por otras máquinas. Asimismo, el único escenario de fallo inyectado probado fue la desconexión del sistema neumático, un modo de fallo común pero no el único posible en el sistema electromecánico completo. Finalmente, la implementación no estuvo exenta de desafíos en la gestión del cronograma. La fase de integración entre la Máquina de Estados Finitos (FSM) y el entorno de simulación Webots presentó una curva de aprendizaje más pronunciada de lo estimado, lo cual se reflejó en una extensión justificada de 3 días en el plan de trabajo original.

\subsection{Reflexión Final Sobre la Transformación Digital de Activos Industriales}
Al término de esta investigación, reflexionamos sobre el impacto de la transformación digital en el tejido industrial latinoamericano. Hemos constatado que la obsolescencia tecnológica no es un destino irreversible, sino un estado transitorio que puede resolverse mediante inteligencia ciberfísica. Nuestro proyecto demuestra que la ecuación de modernización no se limita a la disyuntiva del reemplazo de equipos, sino que se enriquece con una tercera vía: la "digitalización envolvente". Este enfoque permite a las pequeñas y medianas empresas, que conforman la mayor parte del sector agroindustrial, acceder a las ventajas de la Industria 4.0 sin ejecutar inversiones prohibitivas. La combinación estratégica de talento humano en formación de posgrado y herramientas tecnológicas de código abierto tiene el poder de catalizar un cambio profundo en la competitividad del sector productivo.

\section{Recomendaciones}

A la luz de los resultados obtenidos, las lecciones aprendidas y la prospectiva tecnológica, se presentan las siguientes recomendaciones estructuradas para los diferentes actores involucrados, con el fin de maximizar el impacto y la sostenibilidad de esta solución.

\subsection{Recomendaciones para el Laboratorio de Procesos de la Universidad Indoamérica}
Se insta a los responsables del Laboratorio de Procesos a incorporar formalmente el Gemelo Digital desarrollado como un módulo didáctico transversal. Esta herramienta no debe limitarse a una demostración, sino convertirse en un banco de pruebas para asignaturas de automatización industrial, sistemas ciberfísicos y protocolos de comunicación. Se recomienda diseñar guías de laboratorio específicas donde los estudiantes utilicen el SCADA de Node-RED para analizar tendencias de rendimiento y simular fallos, fortaleciendo así sus competencias en diagnóstico de fallos.

Asimismo, se sugiere establecer un protocolo de mantenimiento predictivo básico aprovechando las métricas recolectadas. Aunque el sistema actual no prescribe acciones automáticas, el personal del laboratorio puede utilizar los datos de presión del sistema neumático monitoreado en tiempo real para anticipar ciclos de mantenimiento antes de que ocurra una caída de presión crítica. Para enriquecer el ecosistema de sensórica, se recomienda la adquisición de sensores adicionales de vibración y temperatura que puedan integrarse a los pines de expansión I2C actualmente en uso, mejorando así la riqueza informativa del gemelo sin alterar la arquitectura base.

\subsection{Recomendaciones para la Universidad Indoamérica y la Gestión Curricular}
Desde la perspectiva institucional, se recomienda a la Universidad Indoamérica establecer un laboratorio especializado en Gemelos Digitales y Sistemas Ciberfísicos. La exitosa ejecución de este proyecto de maestría con un presupuesto tan austero establece un precedente para equipar un espacio dedicado con múltiples máquinas \textit{legacy} instrumentadas. Este laboratorio posicionaría a la universidad como un referente en la enseñanza de TIC aplicadas a la industria. Es crucial fomentar proyectos de titulación que sigan esta línea de investigación aplicada, cerrando la brecha entre la teoría académica y las necesidades tangibles de la industria local.

Se recomienda a las autoridades académicas gestionar alianzas estratégicas con el sector agroindustrial de la provincia. Esta transferencia tecnológica directa permitiría calibrar el prototipo en condiciones reales y, al mismo tiempo, ofrecer a los empresarios una ventana de oportunidad para la modernización de bajo riesgo. La inclusión de contenidos sobre protocolos IoT industriales (MQTT, OPC UA) y simulación 3D en la actualización curricular de la Ingeniería en Tecnologías de la Información es igualmente prioritaria para garantizar que los futuros ingenieros dominen estas competencias clave.

\subsection{Recomendaciones Estratégicas para el Sector Agroindustrial}
Para las empresas del sector agroindustrial, este proyecto ofrece un modelo de adopción tecnológica de bajo riesgo. La recomendación principal es realizar un análisis de costo-beneficio que compare el \textit{bypass} ciberfísico presentado aquí con la compra de una máquina nueva equivalente. Nuestros resultados indican que, por menos de 250 dólares, una empresa transita de una operación operativamente "ciega" a un modelo con monitorización en tiempo real, lo cual impacta directamente en los indicadores de disponibilidad.

Se recomienda a los gerentes de planta considerar la monitorización del OEE, incluso sin el factor de calidad automático, como un primer paso para implementar la Mejora Continua. La visibilidad que otorga un tablero SCADA sobre los tiempos de ciclo y los micro-paros permite al supervisor identificar cuellos de botella invisibles a simple vista. La gran ventaja de este enfoque es que se apoya en tecnología \textit{open-source}, lo cual elimina los costos de licenciamiento de software propietario, reduciendo dramáticamente las barreras de entrada para pequeñas y medianas empresas procesadoras.

\subsection{Recomendaciones Técnicas para Iteraciones Futuras del Hardware y Software}
Técnicamente, se recomienda que la próxima iteración del prototipo priorice la incorporación del factor de Calidad en el cálculo del OEE. Esto puede lograrse mediante la integración de una celda de carga y un módulo conversor analógico-digital de alta precisión, añadido al bus I2C existente, para registrar el peso exacto de cada funda. Paralelamente, se sugiere el desarrollo de una aplicación móvil o una interfaz web responsive que actúe como cliente del broker MQTT, permitiendo el monitoreo remoto de la producción desde cualquier dispositivo.

Para mejorar la resiliencia del sistema embebido, se aconseja implementar la lógica de \textit{watchdog} en los ESP32 para reinicios automáticos en caso de bloqueo del firmware. Además, el modelo de la Máquina de Estados Finitos (FSM) actual, que contempla un escenario de falla neumática, debe expandirse para incluir estados correspondientes a fallas eléctricas, atascos de granos o desconexión del bus de comunicación. Finalmente, aunque el consumo es abordable, una auditoría energética del sistema implementado permitiría optimizar los ciclos de sueño profundo de los microcontroladores, una ventaja crucial si el sistema se desplegase en ubicaciones con suministro eléctrico inestable.

\section{Trabajos Futuros}

La culminación de este proyecto de maestría no cierra el ciclo de investigación, sino que abre un abanico de líneas de trabajo futuras que se alinean con la evolución natural de los gemelos digitales y la ingeniería de sistemas ciberfísicos.

\subsection{Evolución hacia un Gemelo Digital Prescriptivo (Nivel 4)}
La línea de investigación más directa que se abre es la migración del actual gemelo de Nivel 3 (Predictivo) a uno de Nivel 4 (Prescriptivo). Esto implica el diseño e integración de algoritmos de inteligencia artificial que no solo anticipen el fallo, sino que recomienden autónomamente la acción correctiva. Se propone explorar el uso de redes neuronales artificiales para modelar el desgaste de los pistones neumáticos basándose en la data histórica de presión recolectada. Del mismo modo, la implementación de algoritmos genéticos podría optimizar en tiempo real la cadencia de los actuadores para maximizar el OEE sujeto a restricciones de fatiga de materiales, transitando de la alerta a la decisión autónoma.

\subsection{Nuevas Líneas de Investigación Aplicada en Ciberfísica}
Se recomienda la derivación de este trabajo hacia la aplicación de la metodología de \textit{bypass} en otras familias de máquinas \textit{legacy} presentes en los talleres de la universidad, como tornos paralelos, fresadoras convencionales o prensas hidráulicas, unificándolas bajo un mismo espacio de monitorización. La integración de tecnologías de Realidad Aumentada (AR) representa un salto cualitativo para el mantenimiento asistido; superponer los datos del gemelo digital sobre la máquina real mediante una Tablet agilizaría las labores de reparación.

Asimismo, la aplicación de tecnologías de registro distribuido (\textit{Blockchain}) podría ser objeto de estudio para establecer una cadena de custodia digital de las variables de operación. Esto permitiría garantizar la inmutabilidad de los registros de producción, un requisito cada vez más demandado en las certificaciones de calidad alimentaria para la exportación de granos. Finalmente, para escenarios donde la latencia de la red sea crítica, se propone investigar la migración de la lógica de control desde la nube hacia un modelo de \textit{Edge Computing}, procesando los datos directamente en los núcleos de los ESP32 para un control de lazo cerrado ultra-rápido.

\subsection{Validación Industrial y Escalabilidad Tecnológica}
El paso fundamental para la continuidad de este proyecto es la validación en un entorno productivo real. Se propone un estudio longitudinal en una planta agroindustrial donde se comparen dos líneas de dosificación idénticas: una operando en modo \textit{legacy} y otra con el gemelo digital implementado. Este estudio permitiría medir de manera exacta el Retorno de la Inversión (ROI) de la modernización. A nivel de conectividad, se recomienda la investigación sobre la sustitución del protocolo MQTT por OPC UA, un estándar con mayores capacidades de seguridad y modelado semántico de datos, preparando la arquitectura para una futura integración horizontal con sistemas MES y ERP. La eventual migración del modelo de simulación de Webots a plataformas de alto fotorrealismo como Nvidia Omniverse se vislumbra como una opción atractiva para entrenar operadores en entornos virtuales ultra realistas antes de que interactúen con el sistema físico modificado.

\begin{verbatim}
=== VERIFICACIÓN DE CITAS ===
NO SE UTILIZARON CITAS BIBLIOGRÁFICAS EXTERNAS EN ESTE CAPÍTULO.
Conforme a las instrucciones de redacción proporcionadas, este capítulo 
sintetiza exclusivamente los hallazgos, métricas y resultados 
obtenidos durante la ejecución del proyecto de titulación 
"Gemelo Digital de Nivel 3 para Máquina Dosificadora de Granos Legacy".
Todos los datos numéricos (OEE, MTTR, latencia, presupuesto) referidos 
en este capítulo provienen de los resultados experimentales propios 
documentados en el Capítulo 3 de esta tesis.
\end{verbatim}
```